\documentclass{article}

\usepackage{url}

\begin{document}

\title{Pairwise Analysis of Correlation Between Genetic Differentiation and Divergence in House Sparrows}
\date{2026-02-17}
\author{Oliver E. Todreas}

\maketitle

\section{Introduction}

To examine the relationship between genetic differentiation and divergence for correlation, those metrics must first be understood. Differentiation, or \textit{fixation index} ($F_{ST}$), is an estimation of population differentiation. Divergence, or in the case of this paper, \textit{per-site divergence} ($d_{XY}$) is the direct nucleotide divergence between populations. They are calculated in fundamentally different ways. $F_{ST}$ is defined as the average frequency of an allele among different populations, balanced against the diversity of a population itself, while $d_{XY}$ is defined simply as the proportion of sites that differ between populations. In other words, one takes into account a populations diversity when doing pairwise comparisons, while the other does not. To determine if these are correlated, we built a pipeline to answer this question. The pipeline consists of filtering, data conversion, parameter and statistics calculations, and plotting using a variety of bioinformatic tools.

\section{Methods}

A shell script was used to filter, index, and convert a VCF file. The filtering step was applied using hard filters in VCFtools at both the genotype and variant level in lieu of a truth set. Only biallelic SNPs passed the filter. Any genotypes with Phred scores below 20 was set to missing, as well as any genotypes with sequencing depths outside the range of 5–70. Then, sites with missingness of 10\% or more were excluded. After filtering, data were indexed using BCFTools. Once indexed, they were converted first to an Intermediate Columnar Format (ICF) directory, then to a Zarr directory using vcf2zarr.

A Python script was used to read the Zarr data into memory, compute $F_{ST}$ and $d_{XY}$, calculate correlation between the two metrics, and plot the results. The filtered Zarr directory was assigned to an Xarray Dataset object using sgkit. A new dimension was merged into the dataset as samples were divided into cohorts based on their populations. The dataset was then divided into two Xarray Datasets, one for each chromosome in the VCF. Windows of 300 variants were then merged into the dataset, and the variant dimension, an Xarray DataSet, was chunked to match the size of the windows, at 1,250 variants per chunk, allowing pairwise $F_{ST}$ and $d_{XY}$ to be calculated and stored in similarity matrices, one for each window of variants using sgkit. Finally, Pearson correlation coefficients were calculated on groups of pairwise $F_{ST}$ and $d_{XY}$ sets for each chromosome using scipy, and data were plotted using matplotlib.

\section{Results \& Conclusions}

Filtering cutoffs were motivated by close examination of the data as well as the method outlined by Dennis et al., who motivated the decision to include only biallelic SNPs (TODO: cite). sgkit was used to generate $F_{ST}$ and $d_{XY}$ because of its flexibility with assigning cohorts and windows, integration of Xarray's strengths like chunking, and because it is the successor to scikit-allel, a popular toolkit that is now only being maintained. Windowing by variant was chosen for its reliability in the number of sites per window, guaranteeing 300 VCF rows of data for each window on which $F_{ST}$ and $d_{XY}$ were computed, except for the final window of each chromosome. A window size of 1,250 variants was chosen because base pairs outnumbered sites in the original data by roughly 40:1, meaning that roughly the same number of windows would be obtained using 1,250-variant windows as there would be 50kb-position windows. Pearson correlation coefficients were used to assess correlation because upon scrutinizing the $F_{ST}-d_{XY}$ plots, many groups appeared right-skewed along $F_{ST}$.

There was more meaningful correlation between $F_{ST}$ and $d_{XY}$ for some populations compared to others (Figure 1). For instance, differentiation and divergence appeared highly correlated between both Lesina and 8N and Lesina and K cohorts on chromosome Z, with $\rho = 0.94$ and $\rho = 0.92$ respectively ($p = 3 \times 10^{-91}$, $p=1\times10^{-80}$). Pearson correlation coefficients were estimated with high certainty for all cohort comparisons and chromosomes ($p<0.05)$, but for the remainder of the cohort pairs, correlation was less clear.

The biological interpretation of these results is that sites on Lesina cohort's Z chromosome that have a high proportion of differing bases compared to the other cohorts also maintain that difference even when diversity within their cohort is taken into consideration, and the same agreement is observed in the lower $F_{ST}$ and $d_{XY}$ values.

\begin{thebibliography}{99}

Dennis TPW, Sulieman JE, Abdin M, Ashine T, Asmamaw Y, Eyasu A, Simma EA, Zemene E, Negash N, Kochora A, Assefa M, Elzack HS, Dagne A, Lukas B, Bulto MG, Enayati A, Nikpoor F, Al-Nazawi AM, Al-Zahrani MH, Khaireh BA, Kayed S, Abdi AA, Allan R, Ashraf F, Pignatelli P, Morris M, Nagi SC, Lucas ER, Hernandez-Koutoucheva A, Doumbe-Belisse P, Epstein A, Brown R, Wilson AL, Reynolds AM, Sherrard-Smith E, Yewhalaw D, Gadisa E, Malik E, Kafy HT, Donnelly MJ, Weetman D. The origin, invasion history and resistance architecture of Anopheles stephensi in Africa. bioRxiv [Preprint]. 2025 Mar 26:2025.03.24.644828. doi: 10.1101/2025.03.24.644828. PMID: 40196515; PMCID: PMC11974716.

\bibitem{sgkit2021}
sgkit Development Team (2021).
\textit{sgkit: Scalable genomics toolkit}.
\url{https://github.com/sgkit-dev/sgkit}

\bibitem{hoyer2017xarray}
Hoyer, S., \& Hamman, J. (2017).
xarray: ND labeled arrays and datasets in Python.
\textit{Journal of Open Research Software}, 5(1), 10.

\bibitem{danecek2021bcftools}
Danecek, P., et al. (2021).
Twelve years of SAMtools and BCFtools.
\textit{GigaScience}, 10(2), giab008.

\end{thebibliography}

\end{document}